% ----- Chapters -----
\chapter{Introduction}

\section{Overview}
Contemporary learning systems are changing from a passive approach to an active one with the assistance of Artificial Intelligence. In the study, we present an Intelligent Tutoring System (ITS) using a Multi-Agent Architecture. The proposed system consists of a multiple agents rather than employing a single chatbot. These agents are designed to simulate a one-on-one tutoring session like a private tutor.

\section{Motivation}
The divide in education technology is no longer just about having access to the internet; it is also about having access to quality, personalized education. E-learning systems have been used traditionally as content management systems but fail to address the Two Sigma Problem \cite{bloom1984}, whereby students receiving individual tutoring outperform those taught in classrooms. We are motivated to provide personalized education by means of autonomous agents that will be able to remember the learner’s past, understand what they upload, and help them when problems arise.

\section{Objectives}
Our main goal is to create a system that adheres to the following objectives:
\begin{itemize}
	\item Contextual Intelligence: Build an LLM Orchestration Layer that manages state across sessions using cache and long-term vector memory.
	\item Autonomous Assessment: Create a pipeline that converts raw media (documents, images, videos) into chunked information and generates evaluative questions based on the document.
	\item Automated Feedback: Create a feedback mechanism that produces reports and learning insights based on existing progress in a topic.
	\item Behavioral Persistence: Design a profiling engine that stores information about the users and their knowledge gaps. This ensures a reduction in the loss of context for the LLM.
\end{itemize}

\section{Scope}
The current study is concentrated on developing a backend system based on the above objectives.

The system will consist of an API Gateway that ensures secure communication,a data layer based on Vector Databases, and an orchestration layer, where different LLMs will fulfill various functions.

\section{SDG Justification}
This project supports United Nations Sustainable Development Goal 4 (Quality Education), SDG 9 (Industry, Innovation and Infrastructure) and SDG 10 (Reduced Inequalities). By using automation for working on time-consuming tasks such as progress tracking, the system offers an affordable, impctful solution for learners worldwide. It tries to balance the inequalities of society in obtaining quality education by providing a personalized learning companion that is available around the clock. Traditional tutoring services can be costly, posing challenges for learners in low-income communities. To alleviate this, the proposed system provides personalized tutoring capabilities that can be accessed anywhere and anytime.

\section{Existing System}
Current Learning Management Systems (LMS) are mostly static. They may not understand the state of mind of the learner, because they cannot perform one-on-one tutoring. MOOC platforms can only imitate classroom-based learning in an online setting.

Thus, current systems struggle to maintain a coherent conversation over prolonged study periods.

\section{Proposed System}
We propose that a Multi-Agent Orchestration system be adopted. In the center of this system, an LLM Orchestrator will direct user inputs towards specific agents according to their intents.

\begin{itemize}
	\item The Content Analysis Agent is concerned with parsing of documents, which are needed for processing of study materials.
	\item The Profiling Agent updates the persistent User Profile Database, accounting for the progress that the learner makes.
	\item The Question Generation Agent aligns the content with the level of difficulty of the user.
    \item The Insight Agent recommends the following actions.
    \item The Proactive Agent interacts with the user via a scheduling capability.
\end{itemize}

\section{Work Plan}
The process of implementation is planned to follow the following stages:

\begin{enumerate}
	\item Stage I: Interagent communication and storage schemas of user data development.
	\item Stage II: The creation of the API Gateway and Orchestrator.
	\item Stage III: Specialized agents' and Vector Database's development and calibration.
	\item Stage IV (Improvement): Implementation of the Proactive Scheduler and testing of the system as a whole.
\end{enumerate}